\section{Conclusiones y trabajos futuros}
\subsection{Conclusiones}
\begin{frame}
  \frametitle{Conclusiones}
  \begin{itemize}
    \item Se han alcanzado los cinco objetivos planteados.
    \item Se observa el \textbf{potencial del método} por aprendizaje para \textbf{automatizar} tareas de metrología monocular.
    \item \textbf{Se dispone de una implementación} de Zhu \emph{et al.}~\footcite{SVMIW}.
      \begin{itemize}
        \item \url{https://github.com/briansenas/TFM_SVMIW}
      \end{itemize}
    \item Se identifican las limitaciones y posibles futuras investigaciones.
  \end{itemize}
\end{frame}

\note{
  Con todo esto, se concluye haber logrado los objetivos del TFM, determinando la posibilidad de
  resolución del problema y su automatización. Permitiendo la extensión del método implementado a más recientes arquitecturas
  de modelos, algoritmos de entrenamiento y conjuntos de datos.
  Aunque la implementación actual no supera al clásico de calibración manual, permitir automatizar y democratizar el proceso en un futuro.
  Los resultados son preliminares, aún así se identifican las diferentes limitaciones y se propone las siguientes posibles líneas de investigación:
}

\subsection{Trabajos futuros}
\begin{frame}
  \frametitle{Trabajos futuros}
  \begin{itemize}
    \item Ajuste de los pesos de las diferentes componentes de la función de pérdida.
    \item Estimación amodal de caja delimitadora~\footcite{kar2015amodalcompletionsizeconstancy}.
    \item Adaptar avances del campo de estimación de profundidad.
      \begin{itemize}
        \item Por ejemplo, utilizando la coherencia entre fotogramas de un video~\footcite{ke2025videodepthvideomodels}.
      \end{itemize}
    \item Mejorar la vectorización del proceso de metrología.
  \end{itemize}
\end{frame}

\note{
  Para evitar el efecto de regresión al a priori un ajuste equilibrado sobre la combinación lineal de las pérdidas,
  junto a un control más metódico sobre el aprendizaje de las capas relacionadas podría regularizar ese comportamiento.
  Adicionalmente, se podría utilizar avances del campo de estimación de profundidad, como la imposición de la coherencia
  entre fotogramas cercanos de un video.
  Por otro lado, cara a ampliar las aplicaciones del método, se podría entrenar a que realice estimaciones amodales
  para tratar con personas u objetos semi-ocultos.
  Esto es todo, gracias! ¿Preguntas?
}
